% §4.9 Smaller groups --- a one-page summary (Q-0312, review flags 4, 12 and 19). Until Q-0312 this
% was section 4.7, with eight statements, and it stood between the step sets of the orbits and the
% theorem on the general linear row, which it delayed (the review's "side branch"). The statements
% moved by paper placement (records/affine/move-q0312-section-4.py; the earlier text is in git):
% prop:similarity-stable-order-families, whose clauses (1)-(2) thm:wedge-axis spends, to the end of
% 4.6; prop:frame-families, prop:fixed-frame-families, prop:similarity-ray-families,
% prop:diagonal-ray-families and prop:diagonal-stable-order-families, with their proofs, to Appendix F
% (appendix-smaller-groups.tex). One linear flow and the similarities with the Loewner order were
% already Appendices G and H. No blueprint node changed; the table rows of the smaller groups stay in
% tab:landscape. No material leaves the paper (the brief: a sequel is the author's decision).
% Round-2 review (Gemini), items 8 and 9, the author's decision D8 (2026-11-03;
% records/affine/reviews/VERIFY-round2-structure-paper.md § 9, option (b)): clauses (3)-(5) of
% prop:similarity-stable-order-families are the node prop:similarity-stable-order-kernels, printed here
% at the end of the paragraph "The stable order" as the subsection's one statement, with the lead
% sentences and the comments that stood beside it in 4.6; clauses (1)-(2) stay at the end of 4.6.

\subsection{Smaller groups}\label{sec:smaller-groups}

A smaller group serves a use that transports fewer shapes. A comparison of an image with its rotated
and rescaled copies, but not with its sheared ones, needs covariance under the similarities only. A
filter applied in the fixed axes of the sampling grid, with a scale for each axis, needs covariance
under the positive diagonal matrices only. Neither group carries every shape to every other, so
neither supports shape selection in the sense of \S\ref{sec:landscape-uses}. Following a shape across
scale and changing it remain, and a smaller group admits at least as many families for them in every
column of Table~\ref{tab:landscape}, strictly more in the cells listed after it. The main gain is in the Loewner column: with the
similarities in place of the general linear group, the classical order keeps non-Gaussian families.
For $0 < \alpha < 2$, $b \ge 0$ and $\kappa > 0$ with $\alpha\kappa^2 - (d+1)\kappa - 1 \ge 0$, one is an
affine Gaussian, trivial when $b = 0$, convolved with a stable kernel whose shape is pulled towards the
isotropic one,
\[
  \mu_\Sigma = N(0, 2b\Sigma) * S_\alpha(\Sigma + \kappa\tr(\Sigma)I) ,
\]
with $N(0, C)$ the centred Gaussian law of covariance $C$. It cascades along every Loewner step and is
covariant under the similarities, and its stable part is rotation invariant only at the multiples of
$I$ (Proposition~\ref{prop:similarity-stable-family}(3)). The statements behind the rows of the two
smaller groups are in Appendices~\ref{app:smaller-groups}--\ref{app:similarity}, except
Proposition~\ref{prop:similarity-stable-order-kernels}, which ends this subsection.

\paragraph{Frames.} The affine image of a product of $d$ cascades on the line cascades along the axes
of its own frame. Only the Gaussian loses its frame, in the sense that its kernel depends on the
aperture only through $A\Sigma A^{\tp}$ and not on the frame $A$ (Proposition~\ref{prop:frame-families}).
Conversely, a family covariant under the positive diagonal matrices that cascades along each axis is
a product of self-decomposable laws on the line, one per axis (Proposition~\ref{prop:fixed-frame-families}).
Such a family gains anisotropy along its axes as fast per step as the Gaussian does and cannot turn,
while the stable orbits turn with a smaller gain per step. Turning alone therefore does not exclude non-Gaussian kernels.

\paragraph{The ray column.} Along rays a family of any group is a continuous covariant field of
self-decomposable laws (Proposition~\ref{prop:gw-ray-families}(2)). For the similarities the orbits of
the shapes form a simplex of dimension $d-1$, and for the positive diagonal matrices they are indexed
by the correlation matrices. In both rows the isotropic member no longer determines the family
(Propositions~\ref{prop:similarity-ray-families} and~\ref{prop:diagonal-ray-families}), and the
similarity row keeps every rotation-invariant self-decomposable law as an isotropic member in every
column, the Loewner column included.

\paragraph{One linear flow.} The families covariant under one group $\{e^{sB}\}$ that cascade along
its curve $e^{sB}\Sigma_0e^{sB^{\tp}}$ lie outside the scheme, the curve not being an order on
$\PSD(d)$. They are the operator-self-decomposable laws with exponent $B$
(Proposition~\ref{prop:one-parameter-families}, Appendix~\ref{app:linear-flow}), and for a contracting
flow a rotation-invariant non-Gaussian one exists exactly when $\tr B < (d+2)\,\lambda_{\min}(\mathrm{sym}\,B)$
(Proposition~\ref{prop:anisotropy-budget}).

\paragraph{The similarities with the Loewner order.} This is the classical order with the smaller
group, and here non-Gaussian families survive. Theorem~\ref{thm:affine-rigidity} needs the whole
general linear group, since its proof transports one rank-deficient step to every other. A kernel is
infinitely divisible, a Gaussian part convolved with a non-Gaussian part of exponent $\psi_\Sigma$
(Proposition~\ref{prop:similarity-cascade-families}). What the order still controls is how far that
part is from rotation invariance, measured by its exponent spread against the eccentricity of the
aperture, both defined in Proposition~\ref{prop:similarity-stable-order-kernels} below: every family has
$\operatorname{es}(\psi_\Sigma) \le 4\operatorname{ecc}(\Sigma)$ (clause (2) of that proposition). The spread
% Round 3 (content review), E section 13: the rank-one limit named.
is bounded uniformly in $\Sigma$ when the rank-one limit of the \Levy{} measures of the part
(Proposition~\ref{prop:similarity-general-sandwich}(4)) stays comparable to the \Levy{} measure of the
isotropic member (Proposition~\ref{prop:similarity-general-bounded-shape}),
and without such a condition it can grow like a power of the eccentricity
(Proposition~\ref{prop:similarity-degenerate-family}). When the profile of the isotropic member is
bounded at the origin, as the \Matern{} profile is, the \Levy{} measure of every part has the same
logarithmic rate at the origin in every direction (Propositions~\ref{prop:similarity-orbit-family}
and~\ref{prop:similarity-general-index-zero}); away from the origin it need not be rotation invariant
(Proposition~\ref{prop:similarity-line-family}).

\paragraph{The stable order.} For the diagonal group the isotropic member must be
operator-self-decomposable for every diagonal exponent of the cone $\mathfrak k_\alpha$, and on the
frame's shapes the cell lies strictly between its neighbours in the diagonal row
(Proposition~\ref{prop:diagonal-stable-order-families}). For the similarities the stable order
$\le_\alpha$ of Proposition~\ref{prop:similarity-stable-order-families} bounds how far a kernel is from
rotation invariance by the \emph{eccentricity} $\operatorname{ecc}(\Sigma)$ of the aperture, the ratio of
its extreme eigenvalues. The main measure is the \emph{exponent spread} of the exponent of the
non-Gaussian part, which is $1$ exactly for a rotation-invariant exponent. The others are condition
numbers of the Gaussian part and of moment matrices of the \Levy{} measure. All are defined in the
statement, and Appendix~\ref{app:similarity} studies them further. The bounds come from a sandwich:
every definite $\Sigma$ lies between two multiples of $I$ in the order
(Proposition~\ref{prop:similarity-stable-order-families}(1)). They are relative to the shape of
$\Sigma$, as they must be, since the stable orbit belongs to the cell.

% shared with blueprint prop:similarity-stable-order-kernels+91a1ab218491
\begin{proposition}[similarity covariance along the stable order: the families and the shape of their kernels]\label{prop:similarity-stable-order-kernels}
Let $d \ge 2$ and $0 < \alpha < 2$, and let $\le_\alpha$ be the stable order on $\PSD(d)$ and
$\varepsilon_\alpha(\Sigma)$ and $t_\alpha(\Sigma)$ its thresholds, as in
Proposition~\ref{prop:similarity-stable-order-families}(1). Write $\mathrm S(d)$ for the similarities $cQ$, $c > 0$,
$Q \in \Orth(d)$, and $\operatorname{ecc}(\Sigma) := \lambda_{\max}(\Sigma)/\lambda_{\min}(\Sigma)$ for
$\Sigma \in \PD(d)$. For a measure $L$ on $\RR^d \setminus \{0\}$ and $s > 0$ put
\[
  \psi_L(\omega) := \int(1 - \cos\langle\omega,x\rangle)\,L(dx) ,
  \qquad
  \mathcal M_s(L) := \int(|x|^2 \wedge s^2)\,\frac{xx^{\tp}}{|x|^2}\,L(dx) ;
\]
for a function $\psi \ge 0$ on $\RR^d$ let the \emph{exponent spread} $\operatorname{es}(\psi)$ be the
infimum of the $K \ge 1$ with $\psi(Q\omega) \le K\,\psi(\omega)$ for every $Q \in \Orth(d)$ and
$\omega \in \RR^d$, the infimum of the empty set being $\infty$; and put
$\operatorname{cond}M := \lambda_{\max}(M)/\lambda_{\min}(M) \in [1,\infty]$ for
$M \in \PSD(d) \setminus \{0\}$. These are the functionals of
Proposition~\ref{prop:similarity-general-bounded-shape}.
\begin{enumerate}
\item \emph{(The kernels of the similarity row.)} Let $G$ be a closed subgroup of $\GL(d)$ with
$\mathrm S(d) \subseteq G$, and $(\mu_{\Sigma,\Sigma'})$ a $(G, \le_\alpha)$-cascade family of kernels.
Every $\mu_{\Sigma,\Sigma'}$ and every kernel from scale zero $\mu_\Sigma$ is infinitely divisible; for
$\Sigma \in \PD(d)$, $\mu_\Sigma$ is a symmetric law with generating triplet $(C(\Sigma), L_\Sigma, 0)$;
the family $(\mu_\Sigma)_{\Sigma \in \PD(d)}$ is covariant under the similarities,
$(\sqrt c\,Q)_*\mu_\Sigma = \mu_{cQ\Sigma Q^{\tp}}$ for $c > 0$, $Q \in \Orth(d)$ and $\Sigma \in \PD(d)$
(so it satisfies the hypotheses of Proposition~\ref{prop:similarity-general-sandwich}); and every
step $\Sigma \le_\alpha \Sigma'$ in $\PD(d)$ has the infinitely divisible increment $\mu_{\Sigma,\Sigma'}$.
For $G = \mathrm S(d)$ the isotropic members $\mu_I$ of these families are exactly the rotation-invariant
self-decomposable laws other than $\delta_0$.
\item \emph{(The sandwich, and the shape.)} In the setting of (1) let $\Sigma \in \PD(d)$,
$\varepsilon := \varepsilon_\alpha(\Sigma)$ and $t := t_\alpha(\Sigma)$, and write
$\psi_\Sigma := \psi_{L_\Sigma}$. Then $C(I) = c_0I$ with $c_0 \ge 0$, and
\[
  \varepsilon c_0I \le C(\Sigma) \le tc_0I , \qquad
  (\sqrt\varepsilon\,I)_*L_I \le L_\Sigma \le (\sqrt t\,I)_*L_I , \qquad
  \psi_I(\sqrt\varepsilon\,\omega) \le \psi_\Sigma(\omega) \le \psi_I(\sqrt t\,\omega) .
\]
Consequently $\operatorname{cond}C(\Sigma) \le \operatorname{ecc}(\Sigma)^{(d+\alpha)/\alpha}$ if $c_0 > 0$,
$\operatorname{cond}\mathcal M_s(L_\Sigma) \le \operatorname{ecc}(\Sigma)^{(d+\alpha)/\alpha}$ for every $s > 0$ if
$L_I \ne 0$, and $\operatorname{es}(\psi_\Sigma) \le 4\operatorname{ecc}(\Sigma)^{(d+\alpha)/\alpha}$. For a
$(G, \text{Loewner})$-cascade family of kernels the same holds with $\varepsilon := \lambda_{\min}(\Sigma)$,
$t := \lambda_{\max}(\Sigma)$ and $\operatorname{ecc}(\Sigma)$ in place of
$\operatorname{ecc}(\Sigma)^{(d+\alpha)/\alpha}$.
\item \emph{(The cell is not the union of its neighbours.)} Let $b > 0$ and $\kappa > 0$, and for
$\Sigma \in \PSD(d)$ put $S_\alpha(\Sigma) := (\Sigma^{1/2})_*S_\alpha(I)$ and
\[
  \mu_\Sigma := N\bigl(0,\, 2b(\Sigma + \kappa\tr(\Sigma)\,I)\bigr) * S_\alpha(\Sigma) .
\]
Then $(\mu_\Sigma)$ is the family of kernels from scale zero of an $(\mathrm S(d), \le_\alpha)$-cascade
family of kernels, and it is not covariant under $\GL(d)$, hence the kernels of no
$(\GL(d), \le_\alpha)$-family (Lemma~\ref{lem:gw-kernels-from-origin}(2)). For $P \in \PSD(d)$ with $0 < \rank P < d$ there
is no positive finite measure $\rho$ with $\mu_{I+P} = \mu_I * \rho$; so $(\mu_\Sigma)$ is the family of
kernels from scale zero of no $(\mathrm S(d), \text{Loewner})$-cascade family.
\end{enumerate}
\end{proposition}

A uniform bound on the anisotropy is impossible in this cell, since the stable orbit belongs to it by
Proposition~\ref{prop:similarity-stable-order-families}(2) with exponent spread
$\operatorname{ecc}(\Sigma)^{\alpha/2}$. The last sentence of clause (2), the bound
$4\operatorname{ecc}(\Sigma)$, holds for every (similarities, Loewner)-cascade family, with no further
hypothesis. The cell contains the restrictions of its two neighbours, the general linear cell of
Proposition~\ref{prop:similarity-stable-order-families}(2) and the (similarities, Loewner) cell
(Lemma~\ref{lem:gw-kernels-from-origin}(1)), and clause (3) gives a family in neither. Whether the
exponent $(d+\alpha)/\alpha$ is sharp is open; the stable orbit has $\alpha/2$. The proof is in
Appendix~\ref{app:proofs}.
